\documentclass[10pt,a4paper]{amsart}
\usepackage[margin=19mm]{geometry}
\usepackage{amsmath,amssymb,mathtools}
\usepackage[scr=rsfs,cal=euler]{mathalfa}
\usepackage{xfrac}
\usepackage[unicode,hidelinks,bookmarks=false]{hyperref}
\newcommand{\ie}{i.e.\ }
\newcommand{\cf}{cf.\ }
\newcommand{\resp}{respectively\ }
\newcommand{\Subsection}{\subsection}
\providecommand{\nobreakdash}{\nobreak\hbox{-}}
\setlength{\emergencystretch}{2em}
\multlinegap0pt
\usepackage{amssymb}
\usepackage[all]{xy}
\usepackage{breqn}
\newtheorem{theorem}{Theorem}
\newcommand{\CC}{{\mathbb{C}}}
\def\H{{\mathcal H}}
\newcommand{\Jac}{\mathrm{Jac}}
\def\L{{\mathcal L}}
\def\O{{\mathcal O}}
\newcommand{\PV}{{\mathrm{PV}}}
\renewcommand{\S}{{\mathcal S}}
\newcommand{\bs}{{\bf s}}
\newcommand{\bx}{{\bf x}}
\newcommand{\p}{{\partial}}
\title{Elements of Saito theory via Batalin--Vilkovisky algebras}
\author{Alexey Basalaev}
\date{Source: arXiv:2601.00132v1 (CC BY 4.0)}
\begin{document}
\maketitle

\noindent\emph{Source and licence.} Elements of Saito theory via Batalin--Vilkovisky algebras. Version arXiv:2601.00132v1 (CC BY 4.0), distributed by arXiv under the Creative Commons Attribution 4.0 International licence, http://creativecommons.org/licenses/by/4.0/, as recorded in the arXiv licence metadata for this exact version and shown on the abstract page https://arxiv.org/abs/2601.00132v1. Full licence notice: \texttt{licenses/arxiv-2601.00132-CC-BY-4.0.txt}.\par\smallskip

\noindent\textit{Editorial scope.} Counted unit: the article's theorem form (source0.tex lines 1120--1127). The article's complete original proof of it is delivered with it (source0.tex lines 1128--1150, one \texttt{proof} environment of 23 lines). No mathematics is written, repaired or re-derived: the statement and the proof are the article's own lines, in the article's own notation, and no retained line is substituted or re-flowed.  Retained uncounted as the source-local context the proof needs: lines 544--549; lines 1092--1094.  Nothing was omitted from any retained span, so the package carries no omission record.  No retained line contains a citation, so the wrapper carries no bibliography and no bibliography entry.  No retained line contains a figure environment, an image-inclusion command or drawing code, so no image is delivered and the package declares no asset dependency.  Editorial formatting only, in addition: an AMS \texttt{amsart} preamble that restates the article's own declarations and macros used by the retained lines, namely the environments \texttt{theorem} on separate counters, together with the standard packages those lines need.  The retained environments print in this document as Theorem 1 in order of appearance, whereas the article prints the same items with the numbering of its own full document.  The complete native source of the article as distributed in the arXiv e-print for this exact version is bundled unchanged as \texttt{sources/191973/source0.tex}.
\begin{align}\label{eq: preBL as cohomology}
    H^\ast(\PV [z], \bar \p_f + z\p) & \cong \H_f^{(0)},
    \\
\notag
    [\phi(x,z) \cdot 1] & \mapsto [\phi(x,z) d^N\bx].
\end{align}

    \begin{equation}\label{eq: main Jzeta}
        J = e^{(F-f)/z}(\zeta).
    \end{equation}

\begin{theorem}\label{theorem: primitive form}
Let $\Phi_f^\omega$ be the $\p$--trivialization associated with the opposite subspace $\L$.
Then the primitive form $\zeta$ associated to $\L$ is found via the recursive formula for its components in the associated good basis:
\begin{align}
    \zeta_{(p)} = - \left[ \pi_{>0} \cdot (\Phi_f^\omega)^{-1} \Big( \sum_{a=1}^p \frac{(F-f)^a}{z^a a!} \zeta_{(p-a)}  \Big) \right]_f, \quad p \ge 1,
\end{align}
where $[-]_f$ stands for taking the classes in $\Jac(f)$.
\end{theorem}

\begin{proof}
    Using the isomorphism $\Upsilon_0$ of Eq.~\ref{eq: preBL as cohomology}, we can equivalently solve Eq.~\eqref{eq: main Jzeta} in $H^\ast(\PV((z^{-1})) \otimes \O_\S,\overline\p_f + z\p)$.
    Slightly abusing the notation, we use the same letters $J$ and $\zeta$ for the respective elements in $H^\ast(\PV((z^{-1})) \otimes \O_\S,\overline\p_f + z\p)$ and $H^\ast(\PV [z] \otimes \O_\S,\overline\p_F + z\p)$.

    Note that under $\Upsilon_0$ and $\Upsilon_1$, the exponential operator $e^{(F-f)/z}$ is mapped to the operator of multiplication by the $\PV^0[[z^{-1}]]$--element $e^{(F-f)/z}$.

    We have by construction $\Phi_f^\omega(J) = J$ because $\Phi_f^\omega$ acts identically on the good basis elements.
    Then Eq.~\eqref{eq: main Jzeta} is equivalent to
    \[
        J = \left[ \pi_{\le 0} \cdot (\Phi_f^\omega)^{-1} \Big( \sum_{k \ge 0} \frac{(F-f)^k}{z^k k!} \zeta \Big)  \right]_f
    \]
    where both $\frac{(F-f)^k}{z^k k!}$ and $\zeta$ inside the brackets are multiplied as $\PV^0$--elements.

    Then we have
    \[
        \left[ \pi_{> 0} \cdot (\Phi_f^\omega)^{-1} \Big( \sum_{k \ge 0} \frac{(F-f)^k}{z^k k!} \zeta \Big) ) \right]_f = 0.
    \]
    Using the linearity of $\Phi_f^\omega$ and the positivity ansatz on $\zeta$, we have
    \[
        (\Phi_f^\omega)^{-1} (\zeta) = \zeta = - \left[ \pi_{>0} \cdot (\Phi_f^\omega)^{-1} \Big( \sum_{k \ge 1} \frac{(F-f)^k}{z^k k!} \zeta \Big) \right]_f
    \]
    because $\Phi_f^\omega$ acts identically on the good basis elements. This concludes the proof because $\Phi_f^\omega$ does not depend on the variables $\bs$, and $F-f \in \CC[\bx] \otimes \CC[\bs]_1$.
\end{proof}

\end{document}
